\chapter{Présentation générale du projet}

\section{Introduction}

Le secteur de la santé génère chaque jour un volume important de données : dossiers des patients, rendez-vous, diagnostics et prescriptions. La qualité du suivi médical dépend en grande partie de la manière dont ces informations sont organisées, conservées et accessibles. Une gestion manuelle ou éparse de ces données conduit rapidement à des pertes d'information, des doublons et des erreurs de planification.

Dans ce contexte, le recours à un système d'information adapté devient indispensable. C'est l'objet de notre projet : concevoir et développer une application web de \textbf{gestion des consultations médicales} qui centralise l'ensemble des données d'un cabinet et facilite le travail des différents intervenants.

Ce chapitre présente le cadre général du projet : le contexte dans lequel il s'inscrit, la problématique à résoudre, les objectifs visés, la solution proposée ainsi que la méthodologie de travail adoptée.

\section{Contexte du projet}

Dans un cabinet médical, trois types d'informations sont au cœur de l'activité quotidienne : les dossiers des \textbf{patients}, le répertoire des \textbf{médecins} et le suivi des \textbf{consultations}. Le dossier du patient regroupe son identité, ses coordonnées et son historique médical ; le répertoire des médecins décrit les praticiens et leurs spécialités ; enfin, chaque consultation relie un patient à un médecin à une date donnée et peut donner lieu à un diagnostic et à une prescription.

Malgré le caractère structuré de ces données, de nombreuses structures continuent de les gérer sur papier ou à l'aide de fichiers dispersés. Cette situation présente des limites évidentes : difficulté à retrouver un dossier, risque de doublons, absence de contrôle sur les créneaux de rendez-vous et suivi médical lacunaire.

Ce projet est réalisé dans le cadre du \textbf{module JEE} (Java Enterprise Edition) de l'École Nationale des Sciences Appliquées d'Agadir. Il vise à mettre en pratique les concepts fondamentaux du développement d'applications d'entreprise en Java : la persistance objet-relationnelle avec \textbf{JPA/Hibernate}, la couche web avec les \textbf{servlets} et les \textbf{JSP}, ainsi que l'organisation du code selon une \textbf{architecture en couches}.

\section{Problématique}

La gestion traditionnelle (papier ou fichiers isolés) des consultations médicales engendre plusieurs difficultés concrètes :

\begin{itemize}
    \item la \textbf{perte ou la détérioration} des dossiers patients au fil du temps ;
    \item les \textbf{doublons} : un même patient peut être enregistré plusieurs fois faute de contrôle d'unicité ;
    \item les \textbf{conflits de planification} : deux consultations peuvent être programmées au même créneau pour un même médecin ;
    \item la \textbf{difficulté de suivi} : l'historique des consultations, des diagnostics et des prescriptions n'est pas centralisé ;
    \item l'absence de \textbf{contrôle d'accès} : aucune distinction claire entre les rôles des utilisateurs.
\end{itemize}

La problématique centrale du projet peut donc être formulée ainsi : \emph{comment concevoir et développer une application web qui centralise la gestion des patients, des médecins et des consultations, tout en garantissant le respect des règles de gestion et un accès sécurisé en fonction des rôles ?}

\section{Objectifs}

L'\textbf{objectif général} du projet est de concevoir et de développer une application web de gestion des consultations médicales, fonctionnelle et structurée, répondant aux besoins d'un cabinet médical.

De cet objectif général découlent des \textbf{objectifs spécifiques} :

\begin{itemize}
    \item gérer les patients : ajout, consultation, modification et suppression de leurs dossiers, avec recherche et contrôle de l'unicité du CIN ;
    \item gérer les médecins : ajout, consultation, modification et suppression, avec recherche et filtrage par spécialité ;
    \item planifier, modifier et annuler les consultations, en évitant les conflits de créneaux et les dates incohérentes ;
    \item permettre au médecin de consulter son planning et son historique, d'établir un diagnostic et une prescription ;
    \item suivre le statut d'une consultation (planifiée, terminée, annulée) ;
    \item sécuriser l'accès à l'application par une authentification et une gestion des rôles ;
    \item offrir un tableau de bord synthétique (nombre de patients, de médecins et consultations du jour).
\end{itemize}

\section{Solution proposée}

Pour répondre à la problématique et atteindre ces objectifs, nous proposons une application web reposant sur les technologies JEE, organisée selon une architecture en couches :

\begin{itemize}
    \item une couche de \textbf{présentation} construite avec des pages JSP enrichies par la bibliothèque JSTL ;
    \item une couche \textbf{contrôleur} composée de servlets qui traitent les requêtes HTTP et orchestrent les actions ;
    \item une couche \textbf{métier} (service) qui implémente la logique applicative et les règles de gestion ;
    \item une couche d'\textbf{accès aux données} (DAO) fondée sur JPA/Hibernate, chargée de la persistance des entités ;
    \item une base de données \textbf{MySQL} qui stocke l'ensemble des informations.
\end{itemize}

La sécurité de l'application est assurée par une authentification par email et mot de passe, complétée par une gestion des rôles (administrateur et médecin) qui limite l'accès aux fonctionnalités. La modélisation du système est réalisée avec le langage \textbf{UML} à travers les diagrammes de cas d'utilisation, de classes et d'états-transitions.

\section{Méthodologie de travail}

Le développement du projet s'est appuyé sur une démarche \textbf{agile}, inspirée du cadre Scrum. Le besoin a d'abord été exprimé sous la forme d'un \textbf{product backlog}, organisé en \emph{epics} (socle technique, gestion des patients, gestion des médecins, gestion des consultations, tableau de bord et authentification) puis décliné en \emph{user stories} de la forme « En tant que\ldots, je veux\ldots, afin de\ldots ».

Chaque user story a été priorisée selon la méthode \textbf{MoSCoW} (Must, Should, Could) et estimée en points de complexité selon la suite de Fibonacci (1, 2, 3, 5, 8). Le travail a ensuite été organisé en \textbf{sprints} successifs, chacun visant un objectif précis, et une \emph{definition of done} a été définie : le code compile et s'exécute sans erreur, et les critères d'acceptation de chaque story sont validés.

Sur le plan technique, les outils suivants ont été utilisés : \textbf{Maven} pour la construction du projet, \textbf{Tomcat} pour le déploiement, \textbf{MySQL} pour la base de données, \textbf{draw.io} pour la réalisation des diagrammes UML et \textbf{Git} pour la gestion de version. Cette démarche itérative a permis de livrer progressivement un socle technique fonctionnel, puis les différents modules métier, tout en contrôlant la qualité à chaque étape.
